% EDITION STATUS: partial; source units 18/18
\documentclass[11pt,a4paper]{article}
\usepackage[no-math]{fontspec}
\setmainfont{Latin Modern Roman}
\setsansfont{Latin Modern Sans}
\setmonofont{Latin Modern Mono}
\usepackage[french,english]{babel}
\usepackage[left=22mm,right=34mm,top=22mm,bottom=22mm,marginparwidth=23mm,marginparsep=5mm,includeheadfoot]{geometry}
\usepackage{amsmath,amssymb}
\usepackage[normalem]{ulem}
\usepackage{xcolor,tikz-cd,marginnote,etoolbox,microtype,graphicx}
\usepackage[unicode,colorlinks=true,linkcolor=LinkInk,urlcolor=LinkInk,bookmarksopen=true]{hyperref}
\definecolor{LinkInk}{HTML}{254E70}
\definecolor{HandInk}{HTML}{254E70}
\definecolor{DeletionInk}{HTML}{8B3A3A}
\definecolor{QuietInk}{HTML}{555555}
\setlength{\parindent}{1em}
\setlength{\parskip}{3pt plus 1pt}
\setlength{\emergencystretch}{1.5em}
\raggedbottom
\normalmarginpar
\frenchspacing
\newif\ifdiplomatic
\newif\ifreading
\newif\ifenglish
\diplomaticfalse
\readingfalse
\englishtrue
\newcommand{\EditionTitle}{A Guide to the Esquisse}
\newcommand{\EditionForm}{Commentary for a mathematical reader}
\hypersetup{pdftitle={Sketch of a Program — Corrected text with translator’s notes},pdfauthor={Alexander Grothendieck},pdfsubject={A scan-checked edition of Esquisse d'un programme, January 1984, Montpellier cote 119}}
\newcommand{\DeclareScan}[4]{%
  \expandafter\def\csname edurl@#1\endcsname{#2}%
  \expandafter\def\csname edkey@#1\endcsname{#3}%
  \expandafter\def\csname edpage@#1\endcsname{#4}}
\def\CurrentScanURL{}
\def\CurrentPageKey{}
\def\CurrentCopy{copy1}
\def\CurrentApparatus{}
\newcounter{appnote}
\newcommand{\EditionPart}[2]{\clearpage\def\CurrentCopy{#1}\section*{#2}}
\newcommand{\scanpage}[1]{%
  \clearpage
  \expandafter\let\expandafter\CurrentScanURL\csname edurl@#1\endcsname
  \expandafter\let\expandafter\CurrentPageKey\csname edkey@#1\endcsname
  \def\CurrentApparatus{}%
  \hypertarget{scan-\CurrentPageKey}{}%
  \typeout{EDITION-SCAN: \CurrentPageKey; PDFPAGE: \the\numexpr\ReadonlyShipoutCounter+1\relax}%
  \ifreading
    \marginnote{\raggedright\scriptsize\sffamily\href{\CurrentScanURL}{119 / \csname edpage@#1\endcsname}}%
  \else
    \noindent{\small\sffamily\href{\CurrentScanURL}{[scan #1]}}\par\smallskip
  \fi}
\newcommand{\mspage}[1]{%
  \ifreading\marginnote{\raggedright\scriptsize\sffamily p.\,#1}[2\baselineskip]%
  \else\noindent{\small\sffamily[p.\,#1]}\par\smallskip\fi}
% A hand macro around a tikz-cd arrow marks that edge only. TikZ defers edge
% drawing, so its explicit color option is required; a surrounding text color
% group alone would not survive until the arrow is drawn.
\ExplSyntaxOn
\NewDocumentCommand{\EditionHand}{m}
  {
    \tl_if_head_eq_meaning:nNTF {#1} \arrow
      {\ifdiplomatic \EditionColorArrow #1 \else #1\fi}
      {\ifdiplomatic \textcolor{HandInk}{#1}\else #1\fi}
  }
\cs_new_protected:Npn \EditionColorArrow \arrow [#1]
  {\arrow[#1,color=HandInk]}
\ExplSyntaxOff
\newcommand{\hand}[1]{\EditionHand{#1}}

\let\over\relax
\newcommand{\over}[2]{\textcolor{QuietInk}{#1}\textcolor{HandInk}{\ensuremath{^{\!\nearrow}}#2}}
\newcommand{\inserted}[2][]{\textcolor{HandInk}{\ensuremath{\langle}#2\ensuremath{\rangle}}}
\newcommand{\deleted}[1]{\ifmmode\textcolor{DeletionInk}{[#1]_{\scriptscriptstyle\mathrm{del}}}\else\textcolor{DeletionInk}{\sout{#1}}\fi}
\renewcommand{\xout}[1]{\ifmmode\textcolor{QuietInk}{[#1]_{\scriptscriptstyle\mathrm{x}}}\else\textcolor{QuietInk}{\sout{#1}}\fi}
\let\TextUline\uline
\let\TextUuline\uuline
\makeatletter
\UL@pixel=0pt
\makeatother
\renewcommand{\uline}[1]{\ifmmode\underline{#1}\else\TextUline{#1}\fi}
\renewcommand{\uuline}[1]{\ifmmode\underline{\underline{#1}}\else\TextUuline{#1}\fi}
\newcommand{\unsure}[1]{\href{\CurrentScanURL}{[#1?]}}
\newcommand{\illegible}[1]{\href{\CurrentScanURL}{[\ifenglish illegible\else illisible\fi\ #1]}}
\newsavebox{\MarginalTestBox}
\newdimen\MarginalLineHeight
\newcount\MarginalLineCount
\newif\ifMarginalFits
\newcommand{\margin}[1]{%
  \begingroup
  \let\MarginalSavedApparatus\CurrentApparatus
  \edef\MarginalSavedAppNumber{\the\value{appnote}}%
  \setbox\MarginalTestBox=\vbox{%
    \hsize=\marginparwidth\linewidth=\marginparwidth
    \parindent=0pt\parskip=0pt\footnotesize
    \global\MarginalLineHeight=\baselineskip
    \raggedright\noindent\textcolor{HandInk}{#1}\par}%
  \directlua{
    local lines = 0
    for n in node.traverse_id(node.id("hlist"), tex.box["MarginalTestBox"].head) do
      lines = lines + 1
    end
    tex.count["MarginalLineCount"] = lines
  }%
  \MarginalFitstrue
  \ifnum\MarginalLineCount>3\relax\MarginalFitsfalse\fi
  \ifdim\dimexpr\ht\MarginalTestBox+\dp\MarginalTestBox\relax>3\MarginalLineHeight\relax\MarginalFitsfalse\fi
  \ifMarginalFits
    \typeout{EDITION-MARGIN: \CurrentPageKey; PLACEMENT: margin; LINES: \the\MarginalLineCount}%
    \marginnote{\usebox{\MarginalTestBox}}%
  \else
    \typeout{EDITION-MARGIN: \CurrentPageKey; PLACEMENT: block; LINES: \the\MarginalLineCount}%
    \setcounter{appnote}{\MarginalSavedAppNumber}%
    \global\let\CurrentApparatus\MarginalSavedApparatus
    \par\smallskip\noindent\begin{minipage}{\linewidth}\hrule\smallskip
    {\footnotesize\textcolor{HandInk}{\ifenglish Marginal note:\else Note marginale :\fi\ #1}}\smallskip\hrule\end{minipage}\par\smallskip
  \fi
  \endgroup}
\newcommand{\authornote}[1]{\margin{#1}}
\newcommand{\diagram}[1]{\par\smallskip{\centering#1\par}\smallskip}
\newenvironment{diagramblock}{\par\begin{center}}{\end{center}\par}
\newcommand{\DeclareVariant}[2]{\expandafter\gdef\csname edvariant@#1\endcsname{#2}}
\newcommand{\AppItem}[2]{%
  \ifcsdef{edvariant@#2}{%
    \ifcsempty{edvariant@#2}{\PackageError{Esquisse}{Empty apparatus entry #2}{Every executed apparatus marker requires its printed note.}}{}%
  }{\PackageError{Esquisse}{Undefined apparatus entry #2}{Every executed apparatus marker requires its printed note.}}%
  \typeout{EDITION-APP-ENTRY: #2; NUMBER: #1}%
  \par\noindent\textsuperscript{#1}\csname edvariant@#2\endcsname}
\newcommand{\EditionApp}[1]{%
  \stepcounter{appnote}%
  \ifmmode{}^{\scriptscriptstyle\theappnote}\else\textsuperscript{\theappnote}\fi
  \begingroup\edef\x{\endgroup\noexpand\gappto\noexpand\CurrentApparatus{\noexpand\AppItem{\theappnote}{#1}}}\x}
\newcommand{\var}[4]{\ifdiplomatic#1\else#2\ifreading\else\EditionApp{#4}\fi\fi}
\newcommand{\scancontext}[2]{\begingroup\def\CurrentScanURL{#1}#2\endgroup}
\newcommand{\copytwo}[2]{\ifdiplomatic\else\ifreading\else\DeclareVariant{#1}{#2}\EditionApp{#1}\fi\fi}
\newcommand{\frapparatus}[1]{\ifreading\else\EditionApp{#1}\fi}
\newcommand{\tnote}[2]{\ifreading\else\DeclareVariant{#1}{\textit{Translator's note.} #2}\EditionApp{#1}\fi}
\newcommand{\finishscan}{\ifx\CurrentApparatus\empty\else\par\bigskip\noindent\rule{.35\linewidth}{.35pt}\par\smallskip{\footnotesize\setlength{\parskip}{3pt}\CurrentApparatus}\fi}
\newcommand{\DeclareNoteReturns}[2]{\expandafter\def\csname edreturns@#1\endcsname{#2}}
\newcommand{\notecall}[3]{\hypertarget{#2}{\hyperlink{note-#1}{#3}}}
\newcommand{\notedef}[2]{\hypertarget{note-#1}{#2}\,{\scriptsize[\csname edreturns@#1\endcsname]}}
\newcommand{\AppReading}[1]{{\def\margin##1{##1}\def\authornote##1{##1}#1}}
\newcommand{\EmptyReading}{\ifenglish\textit{omitted}\else\textit{supprimé}\fi}
% Observed rounded source glyph, cote 119, site 42, PDF page 43.
% Its conventional alphabet assignment remains uncertain; use with \unsure.
% The contour is a vector traced from the reviewed scan detail, not a font alias.
\makeatletter
\newcommand{\sourceglyph}[1]{\mathord{\mathpalette{\EditionSourceGlyph{#1}}{}}}
\newcommand{\EditionSourceGlyph}[3]{%
  \begingroup
  \setbox0=\hbox{$#2 A$}%
  \resizebox{!}{\ht0}{\csname EditionGlyph@#1\endcsname}%
  \endgroup}
\expandafter\newcommand\csname EditionGlyph@119A\endcsname{%
  \begin{tikzpicture}[x=1pt,y=1pt,baseline=0pt]
  \path[fill,even odd rule]
(3.450,10.100) -- (4.950,10.100) -- (5.050,10.000) -- (5.550,10.000) -- (5.650,9.900) -- (5.950,9.900) -- (6.050,9.800) -- (6.250,9.800) -- (6.650,9.600) -- (6.950,9.300) -- (7.050,9.300) -- (7.600,8.750) -- (7.800,8.350) -- (7.800,8.050) -- (7.550,7.800) -- (7.050,7.800) -- (6.250,8.400) -- (6.050,8.500) -- (5.850,8.500) -- (5.350,8.800) -- (3.950,8.800) -- (3.850,8.900) -- (3.650,8.900) -- (3.550,8.800) -- (3.350,8.800) -- (3.150,8.700) -- (2.250,7.900) -- (2.100,7.850) -- (2.100,7.750) -- (1.900,7.550) -- (1.500,6.750) -- (1.500,6.450) -- (1.400,6.350) -- (1.400,5.850) -- (1.300,5.750) -- (1.300,4.650) -- (1.200,4.550) -- (1.200,4.050) -- (1.300,3.950) -- (1.300,2.450) -- (1.400,2.350) -- (1.400,2.150) -- (1.700,1.850) -- (1.900,1.450) -- (1.900,1.250) -- (2.000,1.050) -- (2.150,1.000) -- (2.550,0.600) -- (2.850,0.500) -- (3.100,0.250) -- (3.100,0.150) -- (2.850,0.000) -- (2.250,0.000) -- (2.050,0.100) -- (1.750,0.400) -- (1.650,0.400) -- (0.900,1.150) -- (0.900,1.350) -- (0.700,1.650) -- (0.700,1.850) -- (0.500,2.050) -- (0.300,2.450) -- (0.300,2.750) -- (0.200,2.850) -- (0.200,3.150) -- (0.100,3.250) -- (0.100,3.850) -- (0.000,3.950) -- (0.000,6.650) -- (0.300,7.050) -- (0.400,7.250) -- (0.400,7.450) -- (0.500,7.550) -- (0.500,7.750) -- (0.700,7.950) -- (0.800,8.250) -- (1.950,9.300) -- (2.050,9.300) -- (2.550,9.700) -- (3.150,10.000) -- (3.350,10.000) -- (3.450,10.100) -- cycle
(6.050,5.700) -- (6.250,5.700) -- (6.500,5.450) -- (6.500,5.250) -- (6.600,5.150) -- (6.600,4.650) -- (6.500,4.550) -- (6.500,2.850) -- (6.400,2.750) -- (6.500,2.550) -- (6.500,2.050) -- (6.600,1.950) -- (6.600,1.250) -- (6.950,0.900) -- (7.500,0.750) -- (7.400,0.550) -- (7.400,0.350) -- (7.250,0.100) -- (6.850,0.100) -- (6.250,0.400) -- (5.850,0.400) -- (5.150,0.000) -- (4.750,0.000) -- (4.650,0.100) -- (4.550,0.000) -- (3.650,0.000) -- (3.500,0.150) -- (3.550,0.300) -- (3.750,0.400) -- (4.250,0.400) -- (4.350,0.500) -- (4.850,0.500) -- (4.950,0.600) -- (5.150,0.600) -- (5.500,0.950) -- (5.600,1.150) -- (5.600,1.350) -- (5.700,1.450) -- (5.700,1.950) -- (5.600,2.050) -- (5.600,2.350) -- (5.500,2.450) -- (5.500,2.750) -- (5.400,2.850) -- (5.400,3.350) -- (5.500,3.450) -- (5.500,4.050) -- (5.600,4.150) -- (5.600,5.050) -- (5.800,5.450) -- (6.050,5.700) -- cycle
(7.950,1.900) -- (8.300,1.850) -- (8.300,1.050) -- (8.150,0.900) -- (7.950,0.900) -- (7.850,0.800) -- (7.550,0.800) -- (7.500,0.950) -- (7.600,1.050) -- (7.600,1.450) -- (7.950,1.900) -- cycle;
  \end{tikzpicture}}
\makeatother




\begin{document}
\begin{titlepage}
\centering
\vspace*{18mm}
{\Large Alexander Grothendieck\par}
\vspace{12mm}
{\Huge\bfseries\EditionTitle\par}
\vspace{8mm}
{\Large\EditionForm\par}
\vspace{14mm}
{\large January 1984\par}
\vfill
{\large Université de Montpellier\par fonds Grothendieck, cote 119\par}
\vspace{8mm}
{\normalsize Two typescript copies: sites 2--59 and 78--135\par}
\vspace{8mm}
{\normalsize Edition dated 14 November 2026\par}
\end{titlepage}
\thispagestyle{plain}
\section*{Preface}
\noindent\textbf{Working checkpoint.} This guide currently contains 18 of the 18 requested commentary units. Its scope is recorded with the checked-page coverage.\par\smallskip
This guide follows the sommaire, sections 1–10 and notes (1)–(7) in the order of the typescript. Page references give the original typescript pagination and link to the archive. The guide is based on the checked text and collation; it does not change the text’s mathematics.

\smallskip
AI-transcribed by Claude Opus 5.5 (Anthropic, effort high, Message Batches API); checked against the scans, and this edition built, by GPT-6 Astra (OpenAI, ChatGPT Pro mode); prompts and records by Claude (Anthropic).


The transcription of copy 1 was made with the text layer of the Grothendieck Circle's typeset \textit{Esquisse} (\texttt{EsquisseFr.pdf}) given to the transcribing model as a draft, to compare with after it had transcribed each page; every reading of this edition was decided on the scans.

The published French text and English translation in \textit{Geometric Galois Actions 1}, edited by L. Schneps and P. Lochak (Cambridge University Press, 1997), the Grothendieck Circle's typeset \texttt{EsquisseFr.pdf}, and the Centre for Grothendieckian Studies' draft transcription G94d are named as published texts. They are not reproduced here.

\section*{Rights and credits}
The rights in the original work belong to its author or the author's representatives or heirs. Where we hold rights, in the AI-produced transcription, translation, corrections and typesetting, they are dedicated to the public domain under CC0 1.0. This edition is published for educational, non-profit purposes, and we claim no rights we do not hold. If you hold rights in this work and object to its publication here, please open a GitHub issue and we will take it down.

\clearpage


\typeout{EDITION-APPARATUS-COUNT: \theappnote}
% WORKING COLLATION PREFACE: revise after the explicit comparison reaches copy2 site135.
% This prefatory unit is excluded from the eighteen commentary units.
\section*{Collation note (working)}
The collation accompanying this working guide is still in progress. In the checked portion, copy 2 has oblique pen strokes through the title and author line that remain uncrossed in copy 1 (\href{https://grothendieck.umontpellier.fr/web/viewer.html?file=../119.pdf\#page=80}{copy 2, p.~1, site 79}; \href{https://grothendieck.umontpellier.fr/web/viewer.html?file=../119.pdf\#page=4}{copy 1, p.~1, site 3}). Its deletions of the extra \textit{n} in \textit{inadéquante} and \textit{convenanble} correct typing slips shared by the two copies and therefore enter the corrected copy 1 text under the corresponding shared variant identifiers (\href{https://grothendieck.umontpellier.fr/web/viewer.html?file=../119.pdf\#page=81}{copy 2, p.~2, site 80}; \href{https://grothendieck.umontpellier.fr/web/viewer.html?file=../119.pdf\#page=86}{copy 2, p.~7, site 85}). On p.~8 an oval and connecting line move the typed words \textit{dans ce jeu} to the position after \textit{Par suite,}; this revision is recorded in the copy 2 apparatus while the base text retains copy 1's order (\href{https://grothendieck.umontpellier.fr/web/viewer.html?file=../119.pdf\#page=87}{copy 2, site 86}). The numbered groups in copy 2's upper-left contents margin, its later proof sign and its doubtful star-like signs remain separately described, including their unresolved readings (\href{https://grothendieck.umontpellier.fr/web/viewer.html?file=../119.pdf\#page=79}{sommaire, site 78}; \href{https://grothendieck.umontpellier.fr/web/viewer.html?file=../119.pdf\#page=86}{p.~7, site 85}; \href{https://grothendieck.umontpellier.fr/web/viewer.html?file=../119.pdf\#page=87}{p.~8, site 86}). The handwritten note calls in both contents pages and the two copies' handwritten archive foliations are recorded as shared handwriting, with each observed reading retained; the foliation is not attributed to the author (\href{https://grothendieck.umontpellier.fr/web/viewer.html?file=../119.pdf\#page=3}{copy 1 sommaire, site 2}; \href{https://grothendieck.umontpellier.fr/web/viewer.html?file=../119.pdf\#page=79}{copy 2 sommaire, site 78}). The final account of the complete collation and its counts will replace this working notice after all fifty-eight page pairs have been inspected.


\subsection*{Sommaire}
The contents page arranges the program in ten parts: the \emph{Envoi}; the
``jeu de Légo-Teichmüller'' and the Galois group; number fields attached to a
``dessin d'enfant''; regular polyhedra and finite fields; the passages toward
``topologie modérée'' and ``théories modérées''; ``A la Poursuite des Champs'';
two-dimensional geometry; teaching; and the epilogue. This ordering gives the
reader two practical routes into the mathematical discussion, through the
surfaces and combinatorics of \S\S2--4 and through the proposed foundations of
\S\S5--7, before the return to two-dimensional geometry in \S8. The N.B.
also specifies the document's two systems of notes: asterisks refer to notes
at the bottom of the same page, while the numbered calls refer to the notes
``rajoutées ultérieurement'' collected at the end. In copy 1 the endpoints
1 and 7 of that numbered sequence are handwritten above typed parentheses;
the small dot beside 7 remains marked as doubtful. The French forms retain
the scanned spellings ``astéristiques'' and ``bidimensionelle'', as well as
the accent in ``Légo''. These details belong to the reader's route through
the document and its successive additions.
(\href{https://grothendieck.umontpellier.fr/web/viewer.html?file=../119.pdf\#page=3}{Copy 1, sommaire, site 2}.)



\subsection*{1. Envoi}
The opening locates the program within Grothendieck's application to the
CNRS and explains the work he intends to carry forward. He refers to the
1972 account of his ``Titres et Travaux'', with its list of publications,
and to two boxes of subsequent manuscript notes that have passed through
successive stages of ``décantation''. His account of ``rédaction'' already
qualifies the word: on p.~1 he describes the projected writing as the
further development, clarification, and deepening of ideas, with the
unexpected turns that accompany discovery. The ``Réflexions Mathématiques'',
begun in February 1983 according to this passage, are to make that movement
of thought visible. The sentence crossing pp.~1--2 gives the complementary
operation: an inadequate image is exposed, and the objects previously
encountered askew acquire their mutual coherence. The final paragraph
assigns the ensuing sketch both a retrospective role, covering ten or
twelve years of reflection, and a prospective role as a program for the
coming years. Its stated readers include the Comité National and other
colleagues, former students, and friends. The 1972 report and the manuscript
``Réflexions Mathématiques'' are named here as the author's own points of
reference; the checked passage supplies their role in this application.
\S7 returns to the projected \emph{Réflexions Mathématiques}, while
\S\S9--10 return to the teaching appraisal and the CNRS application.
(\href{https://grothendieck.umontpellier.fr/web/viewer.html?file=../119.pdf\#page=4}{Copy 1, p.~1, site 3};
\href{https://grothendieck.umontpellier.fr/web/viewer.html?file=../119.pdf\#page=5}{p.~2, site 4}.)



\subsection*{2. A ``jeu de Légo-Teichmüller'' and the Galois group}
The second section develops the program from the shared geometric
intuition of surfaces that Grothendieck describes as emerging from his
teaching (pp.~2--3). The transition from combinatorial topology to complex
algebraic curves brings back the fields, rings, and groups named on p.~3,
then the ``multiplicités modulaires'' and their automorphisms. On
pp.~4--5 the family of moduli multiplicities, its gluing and erasing
operations, and the corresponding fundamental groupoids form the
``tour de Teichmüller'', on which he considers the action of the absolute
Galois group. He states the faithfulness of the action on the first
nontrivial level and asks for a description of the tower's action in terms
of that level; p.~5 explicitly says that he has no conjecture yet for a
characterization of the Galois group inside the outer automorphism group.
Note (1), added on p.~52, later withdraws the accompanying judgment
``hors de portée'' and reports a renewed step toward that characterization.
The groupoid language keeps the chosen collections of base points visible
under the symmetries. On pp.~6--7 he formulates the construction principle
in which the first two moduli dimensions supply generators and relations,
sets it beside Demazure's rank principle, and describes the remaining
topological work in terms of the fundamental group at infinity and
``dévissage'' of stratified structures, with an explicit reference to
\S5. On pp.~8--9 the ``petite boite'' and ``grande boite'' turn those
operations into assemblies of surfaces; Thurston's surgery is named,
canonical cellular decompositions point forward to \S3, and real
structures guide the proposed decompositions of moduli multiplicities.
A mathematical slip remains printed on p.~9: the relations
$\rho^3=\sigma^2$, $\rho^4=\sigma^6=1$ imply that
$z=\rho^3=\sigma^2$ satisfies $z^4=z^3=1$, hence $z=1$,
$\rho=1$, and $\sigma^2=1$. The resulting presentation has mutually
inverse maps to the cyclic group of order two, sending $\rho$ to its
identity and $\sigma$ to its generator, and sending that generator back
to $\sigma$. In contrast, the map
$n\longmapsto M_n=\left(\begin{smallmatrix}1&n\\0&1\end{smallmatrix}\right)$
from $\mathbb Z$ to $\mathrm{SL}_2(\mathbb Z)$ is an injective
homomorphism: $M_mM_n=M_{m+n}$, and its upper-right entry recovers $n$.
This establishes the
failure of the displayed presentation to present the group named there.
The readers retain the scanned exponents. For subsequent work on the
Galois-group program, one may name \emph{La Longue Marche à travers la
théorie de Galois} and Drinfeld's
\href{https://www.mathnet.ru/eng/aa199}{\emph{On quasitriangular quasi-Hopf algebras and on a group that is closely connected with}
$\mathrm{Gal}(\overline{\mathbb Q}/\mathbb Q)$}.
(\href{https://grothendieck.umontpellier.fr/web/viewer.html?file=../119.pdf\#page=5}{Copy 1, pp.~2--9, beginning at site 4};
\href{https://grothendieck.umontpellier.fr/web/viewer.html?file=../119.pdf\#page=12}{mathematical slip on p.~9, site 11}.)



\subsection*{3. Number fields associated with a ``dessin d'enfant''}
The third section resumes the chronology of the surface questions and
constructs the passage from a finite map to an algebraic curve and its
arithmetic data. Ladegaillerie's isotopy work enters on pp.~9--10; the
``repères'', ``drapeaux'', or ``biarcs'' and the cartographic groups
$\underline C_2$ and $\underline C_2^+$ then give the dictionary for maps
and sets with an operator group. On p.~11 Grothendieck places the oriented
cartographic group in the fundamental-group description of the sphere
with three punctures, with the order-two condition at the middle
ramification value. The inverse image of $[0,1]$ recovers the map, and
the vertices, edge centers, and face centers become the points above
$0,1,\infty$ (p.~12); the footnote on that page also treats real curves
and unoriented maps. The rational function on p.~13 makes the Galois
action on coefficients explicit. He describes the discovery as a decisive
change in his mathematical center of interest on p.~12, later qualified
by the childhood recollection in note (2), pp.~52--53, and recounts
another change on pp.~14--15: the supposition he and Deligne had found
``dingue'' becomes the result announced by ``Bielyi''. The discussion then
names \emph{La Longue Marche à travers la théorie de Galois}, describes
the work of January--June 1981, and introduces the reconstruction questions
of ``géométrie algébrique anabélienne''. On p.~16 he records a change of
name from ``Théorie de Galois'' to ``théorie de Galois-Teichmüller'', while
returning to the faithful outer action already used in \S2. The p.~17
discussion carries these fundamental groups and their abelianizations
toward the representations associated with Jacobians. Note (3),
pp.~53--54, enlarges the reconstruction program through étale topoi
and localization; note (4), pp.~54--55, develops the ``machine à motifs''
and the provisional priority list. The p.~17 footnote (*)
contains a mathematical slip in the stated general description of the
quotient kernels as abelian products of Tate modules. The puncture-filling
morphism
$\pi_1(\mathbb P^1_{\overline{\mathbb Q}}\setminus\{0,1,\infty\})
\longrightarrow\pi_1(\mathbb P^1_{\overline{\mathbb Q}})=1$
has the entire free profinite group on two generators as kernel. The
assignment of those generators to $(12)$ and $(23)$ gives a continuous
surjection of this kernel onto $S_3$; their products are $(123)$ and
$(132)$ in opposite orders. A product of the abelian groups
$\hat{\mathbb Z}(1)$ cannot have that quotient, since the image of a
commuting pair remains commuting. This concrete morphism locates the
problem in the footnote's kernel description; its wording is preserved
in the readers. The closing paragraph on p.~18 returns to the role of
dreams and visions in mathematical work. The archive catalogue names
\href{https://grothendieck.umontpellier.fr/archives-grothendieck/}{\emph{Cartes. Étude arithmétique} (cote 138) and \emph{La Longue Marche à travers la théorie de Galois} (cotes 140-1--140-4)}.
Among subsequent texts one may
name Shabat and Voevodsky's
\href{https://www.math.ias.edu/vladimir/node/79}{\emph{Drawing Curves over Number Fields}}.
(\href{https://grothendieck.umontpellier.fr/web/viewer.html?file=../119.pdf\#page=12}{Copy 1, pp.~9--18, beginning at site 11};
\href{https://grothendieck.umontpellier.fr/web/viewer.html?file=../119.pdf\#page=20}{kernel description on p.~17, site 19}.)



\subsection*{4. Regular polyhedra and finite fields}
The fourth section starts from the regular maps whose automorphisms act
transitively on frames, translates their oriented version into Galois
coverings by the dictionary of \S3, and includes infinite maps and
tessellations (pp.~18--19). It then follows the geometric realization of
vertices, edges, and faces into affine and projective spaces over arbitrary
base rings. The move to projective space on p.~20 supplies the author's
duality formalism and accommodates centers at infinity. ``Polyèdres
épinglés'', the higher cartographic group, fundamental reflections, and
parameters expressed through doubled cosines lead on p.~21 to the
universal polynomial base ring; specializations over finite fields then
lead back to the arithmetic curves of \S3 (pp.~22--23). A mathematical
slip occurs across pp.~21--22, where the group generated by the fundamental
reflections is identified with the full automorphism group, and the
icosahedron's group is then called simple of order 60. On the usual
centered icosahedron with vertices
$\{(0,\pm1,\pm\varphi),(\pm1,\pm\varphi,0),(\pm\varphi,0,\pm1)\}$,
$\varphi=(1+\sqrt5)/2$, the map $j:v\mapsto-v$ preserves vertices,
distances, and hence all incidences. It is a central involution in its
orthogonal symmetry group. The cyclic coordinate permutation
$(x,y,z)\mapsto(y,z,x)$ also preserves that vertex set and has order three,
so $\langle j\rangle$ is a nontrivial proper central subgroup of the full
group. The determinant morphism to $\{1,-1\}$ selects the
orientation-preserving subgroup as its kernel. The earlier projective
context leaves this calculation intact: on the affine chart of projective
three-space the same map is represented by
$\operatorname{diag}(-1,-1,-1,1)$ and sends $[v:1]$ to $[-v:1]$.
The readers preserve the scanned statement. On pp.~23--24 the
``paradigme combinatoire'' names a further direction suggested by
specialization in exceptional characteristics; the discussion also names
the 27 lines on a cubic surface and the root system $E_7$, which is
retained. Grothendieck records the limits of the work he has actually
pursued, introduces the circumscribed hyperquadric, and on p.~25 proposes
a modern rewriting of Klein's book as a task for a young researcher.
Note (5), on p.~55, supplements the account of preliminary work with
the arithmetic action on congruence maps and the ``bi-icosaèdre''.
Coxeter's books and Klein's book are the works explicitly invoked here;
the geometric-combinatorial viewpoint continues into the foundational
questions of \S5.
(\href{https://grothendieck.umontpellier.fr/web/viewer.html?file=../119.pdf\#page=21}{Copy 1, pp.~18--25, beginning at site 20};
\href{https://grothendieck.umontpellier.fr/web/viewer.html?file=../119.pdf\#page=24}{reflection-group wording on pp.~21--22, site 23};
\href{https://grothendieck.umontpellier.fr/web/viewer.html?file=../119.pdf\#page=25}{site 24}.)



\subsection*{5. A hue and cry against so-called ``general'' topology, and heuristic reflections toward a so-called ``tame'' topology}
The fifth section starts on p.~25 with the problem of reconstructing a
stratified structure from canonical ``pièces de construction'', and
retraces a change from the successive singular loci of a variety to the
stratifications arising from degeneration in moduli problems. The
Mumford-Deligne compactifications on p.~26 connect this problem with the
Teichmüller tower of \S2. The two-stratum model uses a tubular neighborhood
$T$ of $Y$ in $X$, its complement $U=X\setminus\mathrm{Int}(T)$, and a
gluing map $f:\partial U\longrightarrow Y$; p.~27 reconstructs the space
by the stated amalgamated sum. The description on p.~26 contains a
mathematical slip in the base of the associated cone bundle: the preceding
bundle has projection $p:\partial T\longrightarrow Y$, and its fiberwise
cone is explicitly
$((\partial T\times[0,1])\amalg Y)/((z,0)\sim p(z))$,
with projection $[z,t]\mapsto p(z)$ and $y\mapsto y$ to $Y$.
This formula exhibits the base of that construction; the printed phrase
``fibré en cônes sur $T$'' remains unchanged in the readers. On pp.~27--28
the lack of canonical choices leads to ``catégories isotopiques'',
resuming Ladegaillerie's work from \S3, and then to the author's account
of the limits of his initial hopes in the presence of wild phenomena.
Hironaka's work prompts the renewed search for a ``topologie modérée''
(pp.~28--29). The polemic on pp.~29--30 explains the geometric
requirements behind that search. From p.~30 onward Grothendieck prefers
an axiomatic treatment encompassing several tame theories, with
comparison maps between spaces of isomorphisms and their homotopy types
(p.~31); triangulability is the strongest axiom in his provisional system,
and he explicitly records the unproved state of the proposed comparison
theorem (p.~32). The discussion of fields and point dimensions on
pp.~32--33 explains one advantage he sees in retaining several theories.
Another mathematical slip occurs in the collar description on p.~33,
where the product is written with $T$ as its factor. The elementary collar
of $T=[-1,1]$ is the map
$c:\{-1,1\}\times[0,\varepsilon)\longrightarrow[-1,1]$,
$c(s,t)=s(1-t)$, with $0<\varepsilon<1$; its inverse on the two endpoint
neighborhoods is $x\mapsto(\operatorname{sgn}(x),1-|x|)$.
The displayed isomorphism uses $\partial T$ as the product factor;
the source's plain $T$ is preserved. On p.~34 he allows that the proposed
equivalence of isotopy categories may be too optimistic, while defending
the content of the dévissage already developed. Local retractions provide
additional structure and lead to spaces indexed by flags in the ordered
set of strata, but he immediately notes that these extra data are absent
from the geometric situations that motivated the work. On p.~35 the
moduli example replaces the index set by the category of ``graphes MD'',
and the argument returns to the combinatorial tower of \S2. The
``topossique'' test on pp.~35--36 carries the same heuristic into étale
topoi and the Galois-compatible dévissage of profinite Teichmüller groups;
Grothendieck proposes a ``voisinage tubulaire canonique d'un sous-topos''
and a comparison with the transcendental setting. The final change of
emphasis, on p.~36, is that direct geometric arguments can probably
short-circuit much of this theory for the moduli applications, although
he retains its value as a guide. The later note (6), pp.~55--56,
qualifies the polemic by recognizing the place and beauty of wild
phenomena; note (7), pp.~56--57, gives the morphism data and the hexagonal
diagram behind the flag-indexed construction. Related titles are
\href{https://grothendieck.umontpellier.fr/archives-grothendieck/}{\emph{Topologie modérée} (cotes 120--121), \emph{Espaces stratifiés et voisinages côniques} (cote 150), and \emph{Topos stratifié} (cote 151)},
and Lou van den Dries's
\href{https://assets.cambridge.org/97805215/98385/frontmatter/9780521598385_frontmatter.pdf}{\emph{Tame Topology and O-minimal Structures}}.
(\href{https://grothendieck.umontpellier.fr/web/viewer.html?file=../119.pdf\#page=28}{Copy 1, pp.~25--36, beginning at site 27};
\href{https://grothendieck.umontpellier.fr/web/viewer.html?file=../119.pdf\#page=29}{cone-bundle passage on p.~26, site 28};
\href{https://grothendieck.umontpellier.fr/web/viewer.html?file=../119.pdf\#page=36}{collar passage on p.~33, site 35};
\href{https://grothendieck.umontpellier.fr/web/viewer.html?file=../119.pdf\#page=39}{closing reassessment on p.~36, site 38}.)


\subsection*{6. ``Differentiable theories'' (in the manner of Nash) and ``tame theories''}
The sixth section first continues the foundational program of \S5:
p.~37 proposes a dévissage of a proper tame map, simultaneous with a
finite family of closed subspaces, and introduces ``espaces globalement
modérés'' through complements in compact tame spaces. The specified
comparison of compactifications is the map
$g:X'^{\wedge}\longrightarrow X^{\wedge}$ restricting to the stated
bijection over $X$; p.~38 uses scheme compactifications in the examples
and extends the proposed framework to maps without the initial properness
requirement. There is an index slip on p.~37: the filtration is explicitly
decreasing, so its inclusion maps
$Y^i\hookrightarrow Y^{i-1}$ send every point of $Y^i$ into
$Y^{i-1}$. Consequently the printed ``open stratum''
$Y^i\setminus Y^{i-1}$ is empty. The readers preserve the exponent
$i-1$. After reporting what he had learned from Mebkhout, Grothendieck
states the ``principe des choix anodins'' on pp.~38--39: his criterion
examines the homotopy types of the spaces in which arbitrary construction
choices are made. This leads from isotopy categories, retaining the
components of isomorphism spaces, to the full homotopy types in the
``$\infty$-isotopique'' viewpoint. The surface examples return to the
anabelian cases and to Ladegaillerie's thesis from \S3; the discussion of
1-isotopy categories gives the corresponding bicategorical language.
At the end of p.~39 and the beginning of p.~40 he explains why he has
argued at such length for foundational work he does not intend to make
his own immediate priority. He then turns back to ideas from the late
1960s inspired by Nash, replacing the axioms for tame subsets by axioms
for rings of germs and the resulting smooth manifolds. The rounded
ring-system symbol on pp.~40--42 has an unresolved alphabet assignment;
the readers reproduce its checked contour and keep the doubt marks.
The stability requirements successively concern differential constructions,
the Weierstrass preparation condition and the footnote's henselian
formulation (p.~40), analytic continuation, and integration of Pfaff
systems (p.~41). The refusal to require algebras over all of
$\mathbb R$ leaves room for a smaller field of constants, pointwise
residue fields, and a hierarchy of ``fonctions élémentaires'' and
``constantes élémentaires''; these notions return to the fields discussed
in \S5, pp.~32--33. On p.~42 he judges the interesting questions about
the relation between a finitely generated theory and its field of
rationality ``out of reach'', a judgment explicitly withdrawn in note (1)
on p.~52, then explains how the older differentiable
ideas have become relevant again through the tame program. The final
proposal associates a tame theory with a differentiable one and studies
the differentiable structure off a closed nowhere dense tame subset,
with agreement of two such structures away from a further nowhere dense
subset. The closing sentence presents this as a setting for the more
precise comparison questions raised by the stratified dévissage of \S5.
(\href{https://grothendieck.umontpellier.fr/web/viewer.html?file=../119.pdf\#page=40}{Copy 1, pp.~37--42, beginning at site 39};
\href{https://grothendieck.umontpellier.fr/web/viewer.html?file=../119.pdf\#page=41}{principle of choices on pp.~38--39, beginning at site 40};
\href{https://grothendieck.umontpellier.fr/web/viewer.html?file=../119.pdf\#page=43}{Nash-inspired axioms on pp.~40--42, beginning at site 42};
\href{https://grothendieck.umontpellier.fr/web/viewer.html?file=../119.pdf\#page=45}{closing comparison on p.~42, site 44}.)


\subsection*{7. Pursuing Stacks}
The seventh section opens on p.~42 with the work begun in March of the
previous year on the ``fondements de l'algèbre (co)homologique non
commutative'', also called ``algèbre homotopique''. On p.~43 Grothendieck
describes the intended two-volume \emph{A la Poursuite des Champs}, its
successive expansions, and its place at the beginning of
\emph{Réflexions Mathématiques}, already announced in \S1. His assessment
first places this work somewhat at the margin of the newer interests
described in the preceding sections, as a return to questions from before
1970. He names Verdier, Giraud, and Illusie, acknowledges Quillen's model
categories, and on p.~44 explains the internal structures and
noncommutative coefficients that remain the focus of his own inquiry.
The correspondence with Breen in 1975 leads into the proposed role of
$\infty$-groupoids as models of homotopy types, with truncated types for
n-groupoids; the subsequent correspondence with Brown brings him back
to this problem. The account records two changes on p.~44: a broader
notion initially called an $\infty$-groupoid is renamed ``champ en
$\infty$-groupoïdes'', or ``champ'' over the punctual topos, and
\emph{Histoire de Modèles} becomes a large, unforeseen digression before
the intended treatment of stacks. The relation to \S6 is through these
homotopy types and higher groupoids; p.~45 also explicitly reconnects the
work with the moduli multiplicities and their stratifications from
\S\S2 and 5, through a higher-dimensional van Kampen question. The same
page returns to the 1975--76 reflections on a ``complexe de De Rham à
puissances divisées'' and to ``schématisation'' of homotopy types, while
qualifying the corresponding chapter as particularly hypothetical.
The final paragraph revises the image of merely paying a debt to the
1955--70 period: unforeseen developments, represented by the nested tales
of the \emph{Thousand and One Nights}, account for the present work's
unexpected scale. The archive catalogue names
\href{https://grothendieck.umontpellier.fr/archives-grothendieck/}{\emph{Pursuing stacks. The Modelizing Story} (cotes 134-1--134-8) and \emph{Complexe de De Rham à puissance divisée} (cote 136)};
among later literature one may name Denis-Charles Cisinski's
\href{https://smf.emath.fr/publications/les-prefaisceaux-comme-modeles-des-types-dhomotopie}{\emph{Les préfaisceaux comme modèles des types d'homotopie}}.
(\href{https://grothendieck.umontpellier.fr/web/viewer.html?file=../119.pdf\#page=45}{Copy 1, pp.~42--45, beginning at site 44};
\href{https://grothendieck.umontpellier.fr/web/viewer.html?file=../119.pdf\#page=47}{renaming and digression on p.~44, site 46};
\href{https://grothendieck.umontpellier.fr/web/viewer.html?file=../119.pdf\#page=48}{connections and closing reassessment on p.~45, site 47}.)


\subsection*{8. Digressions in two-dimensional geometry}
The eighth section, beginning on p.~45, gathers the more concrete
two-dimensional questions arising from teaching and research supervision.
On p.~46 Grothendieck describes a repeated sequence: a promising initial
exploration fails to become a shared undertaking with students, and
stronger interests then displace it. Plane polygons and their moduli
varieties form the first theme, including the unexpected appearance of
algebraic geometry. The next theme concerns immersed curves and surfaces
with boundary. The ``stable'' and ``cellulaire'' conditions lead back to
the maps of \S3, while the ``balayage'' of a branch across a double point
and the erasure or creation of a ``bigône'' supply the elementary
operations for a dynamic study (pp.~46--47). The stated questions concern
classes of immersions and the extensions of a prescribed boundary
immersion; the reported tangible result is a canonical dévissage into
indecomposable immersions by ``téléscopage''. His assessment changes with
the third theme on p.~47: work on systems of lines and pseudolines has
brought more productive exchanges with students, although their joint
work has remained elementary. On p.~48 he describes the recent
combinatorial construction of a ``surface modulaire'' parametrizing an
additional pseudoline's relative positions, or the system's
``affinisations'' through a choice of pseudoline at infinity. The
catalogue names the related archive dossier
\href{https://grothendieck.umontpellier.fr/archives-grothendieck/}{\emph{Système de pseudo-droites} (cote 154)}.
The final movement returns explicitly to \S3: since 1977, the conformal
realization of maps, their algebraic curves in the oriented compact case,
and the associated metric possibilities have remained behind these
combinatorial questions. Thurston's geometric ideas reinforce that
expectation. At the transition to p.~49, however, the account of his
students' difficulties leads to the judgment that teaching has often
burdened rather than stimulated his geometric reflection; this judgment
opens the broader appraisal in \S9.
(\href{https://grothendieck.umontpellier.fr/web/viewer.html?file=../119.pdf\#page=48}{Copy 1, pp.~45--49, beginning at site 47};
\href{https://grothendieck.umontpellier.fr/web/viewer.html?file=../119.pdf\#page=49}{elementary operations on pp.~46--47, beginning at site 48};
\href{https://grothendieck.umontpellier.fr/web/viewer.html?file=../119.pdf\#page=51}{moduli surface and conformal viewpoint on p.~48, site 50};
\href{https://grothendieck.umontpellier.fr/web/viewer.html?file=../119.pdf\#page=52}{transition to the teaching appraisal on p.~49, site 51}.)


\subsection*{9. An appraisal of teaching activity}
The ninth section returns on p.~49 to the university teaching begun in
1970 and to the change of professional milieu described at the beginning
of the report. Grothendieck recognizes what this experience taught him
and its part in renewing his mathematical outlook, then calls his
research-level teaching a failure by the criterion of creating a place
where something happens and something ``passe'' among people engaged in
a common adventure. The two exceptions he records, around 1974--76, are
Ladegaillerie's isotopy work already discussed in \S3, and
Contou-Carrère's work on local and global Jacobians over general base
schemes. The rest of the appraisal concerns teaching the rudiments of
the mathematician's craft to some twenty research students. Its final
paragraph on p.~50 qualifies the negative account by recalling work
better than merely acceptable, sometimes pursued beyond a single
project and carried through with taste. He nevertheless anticipates few
opportunities for these students to practice that craft professionally.
This reassessment develops the disappointment described at the end of
\S8 and supplies the immediate background to his application for
secondment in \S10; it is the author's account at the date of the report.
(\href{https://grothendieck.umontpellier.fr/web/viewer.html?file=../119.pdf\#page=52}{Copy 1, pp.~49--50, beginning at site 51};
\href{https://grothendieck.umontpellier.fr/web/viewer.html?file=../119.pdf\#page=53}{closing paragraph on p.~50, site 52}.)


\subsection*{10. Epilogue}
The epilogue on pp.~50--51 connects the program to the CNRS application
announced in \S1. Following the teaching appraisal of \S9, Grothendieck
states that this university activity has ceased to be useful to him or
to others, and explains his request for secondment, the competing claims
on scarce research posts, and the university position that his departure
would release. The application itself has prompted the present sketch;
the table of contents was added afterward to help readers, while the
remarks about ten or a hundred years of work measure the scope he
perceives in its themes. On p.~51 he states their unity through the
geometry of surfaces, bringing together the tower and maps of \S\S2--3,
the geometric constructions of \S4, and the later foundational and
teaching discussions. His final retrospective change is a renewal of
approach to ``motifs'', whose fascination he says has persisted from his
earlier mathematical work; the footnote directs the reader to the
``divagations motiviques'' at the end of the 1972 \emph{Esquisse
Thématique}, pp.~17--18. The archive catalogue also names
\href{https://grothendieck.umontpellier.fr/archives-grothendieck/}{\emph{Motifs et théorie de Hodge} (cote 18)}.
The closing sentences describe his release from the self-imposed burden
of endless ``tâches'' and the greater lightness he associates with age.
The main report ends with the date January 1984; the subsequently added
numbered notes follow it.
(\href{https://grothendieck.umontpellier.fr/web/viewer.html?file=../119.pdf\#page=53}{Copy 1, pp.~50--51, beginning at site 52};
\href{https://grothendieck.umontpellier.fr/web/viewer.html?file=../119.pdf\#page=54}{motives, retrospective footnote, and closing date on p.~51, site 53}.)


\subsection*{Note (1)}
The first later note, on p.~52, explicitly retracts the judgment ``hors
de portée'' used in \S2, p.~5, and again in \S6, p.~42. Grothendieck now
calls the expression hasty and unfounded, treating it as a report on the
speaker's ideas and engagement rather than a measure of a problem's
mathematical possibilities. The joke about La Palisse accompanies his
references to the conjectures of Mordell, Tate, and Chafarévitch. He then
describes a change occurring in the days after he wrote the main report:
renewed contact with the earlier questions has brought a further property
of the outer Galois action on the fundamental group of an algebraic curve
to his attention. He presents this as progress toward the algebraic
characterization sought in \S2 and places it, together with the
``conjecture fondamentale'' of \S3, at the center of his anabelian
interests. The note gives neither that new property's statement nor a
proof; its function here is to record the reassessment and the renewed
direction of work.
(\href{https://grothendieck.umontpellier.fr/web/viewer.html?file=../119.pdf\#page=55}{Copy 1, p.~52, site 54};
\href{https://grothendieck.umontpellier.fr/web/viewer.html?file=../119.pdf\#page=8}{first call on p.~5, site 7};
\href{https://grothendieck.umontpellier.fr/web/viewer.html?file=../119.pdf\#page=45}{second call on p.~42, site 44}.)


\subsection*{Note (2)}
The second note, beginning on p.~52 and ending on p.~53, qualifies the
claim in \S3, p.~12, that no mathematical fact had made an impression
comparable to the discovery attached to a child's drawing. Grothendieck
recalls learning the definition of a circle from Maria, an internee
who gave him lessons at Rieucros when he was about twelve. The definition's
simplicity and obviousness transformed what had seemed the mysterious
``rotondité parfaite'' of the circle. He identifies this retrospectively
as his first glimpse of the creative power of a good definition, a
``formulation qui décrit l'essence'', and says that its fascination has
persisted. The note thus supplies the explicit earlier exception to the
turning point narrated in \S3 while returning to the role of finding a
formulation that the opening discussion of writing in \S1 had already
emphasized.
(\href{https://grothendieck.umontpellier.fr/web/viewer.html?file=../119.pdf\#page=55}{Copy 1, pp.~52--53, beginning at site 54};
\href{https://grothendieck.umontpellier.fr/web/viewer.html?file=../119.pdf\#page=56}{ending at site 55};
\href{https://grothendieck.umontpellier.fr/web/viewer.html?file=../119.pdf\#page=15}{call in \S3, p.~12, site 14}.)


\subsection*{Note (3)}
The third note, on pp.~53--54, extends the program attached to \S3,
p.~16, beyond the varieties called anabelian. Grothendieck proposes a
description of the category of schemes of finite type over the absolute
base through systems of profinite groups expressing their étale topoi,
and then envisages passage to more general schemes by limits. His phrase
``qui fait semblant'' marks the presentation's deliberate avoidance of
rings and equations while retaining the schemes' étale topoi. The
explicit qualification, introduced as ``un grain de sel'', is that the
reconstruction is to concern the category obtained by making universal
homeomorphisms invertible; in the finite-type setting under discussion
he describes these morphisms as finite, radicial, and surjective. This
places the proposed dictionary of \S3 in the language of categories and
localization and supplies further context for the priorities listed in
note (4). The closing paragraph on p.~54 presents this ambition as an
outcome of Galois theory and answers the familiar complaints about its
purpose and generality. The note enlarges and qualifies the proposed
reconstruction; it does not supply a completed construction of that
dictionary. The related archive title \emph{La Longue Marche à travers
la théorie de Galois} is already named with \S3.
(\href{https://grothendieck.umontpellier.fr/web/viewer.html?file=../119.pdf\#page=56}{Copy 1, pp.~53--54, beginning at site 55};
\href{https://grothendieck.umontpellier.fr/web/viewer.html?file=../119.pdf\#page=57}{closing paragraph at site 56};
\href{https://grothendieck.umontpellier.fr/web/viewer.html?file=../119.pdf\#page=19}{call in \S3, p.~16, site 18}.)


\subsection*{Note (4)}
The fourth note, on pp.~54--55, expands the motivic direction attached to
\S3, p.~17. Grothendieck describes the arithmetic structure on
$S\ell(2,\mathbb Z)$ as a ``machine à motifs'', recalls experiments begun
in 1981, and specifies the input homomorphisms into $G_o(\mathbb Z)$ and
the group-scheme extensions used for the resulting Galois
representations. The tautological case leads in his account to an
extension of $G\ell(2)_{\mathbb Z}$ by a two-dimensional torus and to the
representations associated with the two named complex multiplication
fields. He records that this construction principle has not yet been
followed through. The passage from that unfinished exploration to the
five-part priority list on p.~55 makes the program's current order
explicit while calling the order provisional: construction of the
Teichmüller tower, its profinite automorphism group and the Galois
subgroup, the motivic machine, the anabelian dictionary and fundamental
conjecture, and Fermat's problem. Thus items a) and b) return to \S2,
item c) to the Jacobian representations of \S3, and item d) to the broader
reconstruction of note (3). The parenthesis saying that the fundamental
conjecture may be less ``hors de portée'' than supposed continues the
explicit change of mind recorded in note (1). The ``paradigme profini'',
inertia subgroups, positive characteristic, and the absolute ring are
named here as parts of the intended expansion, without a proposed full
formulation being supplied.
(\href{https://grothendieck.umontpellier.fr/web/viewer.html?file=../119.pdf\#page=57}{Copy 1, pp.~54--55, beginning at site 56};
\href{https://grothendieck.umontpellier.fr/web/viewer.html?file=../119.pdf\#page=58}{priority list on p.~55, site 57};
\href{https://grothendieck.umontpellier.fr/web/viewer.html?file=../119.pdf\#page=20}{call in \S3, p.~17, site 19}.)


\subsection*{Note (5)}
The fifth note, on p.~55, adds an exception to the account in \S4,
p.~24, of preliminary work and geometric housekeeping. It identifies
the more delicate combinatorial interpretation of regular maps attached
to congruence subgroups and the effort to express their arithmetic
Galois action through the cyclotomic character. The ``cartes de
congruence'' thus return to both the regular maps of \S4 and the
cartographic arithmetic of \S3. The concrete change of perspective is
the one recorded in the final sentence: the combinatorial
``bi-icosaèdre'', first developed in a C4 course for geometric reasons,
subsequently provided a way to express the action on the category of
icosahedral maps, called here ``cartes de congruence d'indice 5''.
That wording is retained. The note also recalls, without developing it
further, the earlier work on cubic complexes; its emphasis is the
arithmetic use that emerged from a teaching construction.
(\href{https://grothendieck.umontpellier.fr/web/viewer.html?file=../119.pdf\#page=58}{Copy 1, p.~55, site 57};
\href{https://grothendieck.umontpellier.fr/web/viewer.html?file=../119.pdf\#page=27}{call in \S4, p.~24, site 26}.)


\subsection*{Note (6)}
The sixth note, beginning on p.~55 and ending on p.~56, qualifies the
tame-topology discussion attached to \S5, p.~31. Grothendieck identifies
the isomorphism classes of compact tame spaces with those of the
piecewise linear theory, while recalling the stability requirements that
motivate his earlier objections to the latter as a complete context for
the program. He presents this as a rehabilitation, in the intended
setting, of the ``Hauptvermutung''. The important change of emphasis
comes on p.~56: the polemic against wild phenomena as obstacles in the
chosen foundations is accompanied by an explicit recognition of their
reason for being and their own beauty. The proposed foundations are to
place them as ``cas limites'' rather than make them disappear. This
reassessment belongs beside the construction and comparison requirements
of \S5, pp.~27--33, and the treatment of full homotopy types in \S6,
pp.~38--39. The note states the intended comparison and its interpretation;
the comparison's proof is not developed here.
(\href{https://grothendieck.umontpellier.fr/web/viewer.html?file=../119.pdf\#page=58}{Copy 1, pp.~55--56, beginning at site 57};
\href{https://grothendieck.umontpellier.fr/web/viewer.html?file=../119.pdf\#page=59}{closing qualification on p.~56, site 58};
\href{https://grothendieck.umontpellier.fr/web/viewer.html?file=../119.pdf\#page=34}{call in \S5, p.~31, site 33}.)


\subsection*{Note (7)}
The final note, on pp.~56--57, supplies the data behind the flag-indexed
system of spaces mentioned in \S5, p.~34. Its ``strates déployées'',
``tubes de raccord'', ``fibres de raccord'', and ``espaces de jonction''
describe the unfolding of the stratified structure. A target-index slip
occurs in the prose on p.~56: both displayed morphisms there are written
$X_{i,j}\longrightarrow X_i$, the first described as a bordering inclusion
and the second as a proper fibration. The diagram and its stated legend
instead give the hooked immersion
$a_{ij}:X_{ij}\hookrightarrow X_j$ and the ordinary proper fibration
$b_{ij}:X_{ij}\longrightarrow X_i$. Thus it is the target of the first
prose morphism that conflicts with the diagram; both printed prose
targets remain in the readers. The lower right cartesian square identifies
$X_{ijk}$ with $X_{ij}\times_{X_j}X_{jk}$ through its two projections,
using $a_{ij}$ and $b_{jk}$ as the maps to $X_j$. The author's NB stresses
that the data still include the morphism
$X_{ijk}\longrightarrow X_{ik}$, since the immersion
$X_{ik}\longrightarrow X_k$ is not required to be an embedding. On p.~57
he recasts the unfolding of a stratified tame space using the sum $X_*$
of its unfolded strata and the order relation with graph $X_{**}$:
$pr_1:X_{**}\longrightarrow X_*$ is the proper locally trivial fibration,
and $pr_2:X_{**}\longrightarrow X_*$ the bordering embedding. For tame
multiplicities the account uses a category structure in place of that
ordered-object structure, with the maps
$X_{ijk}\longrightarrow X_{ik}$ furnishing composition. The note thus
turns the extra structures and reconstruction problem of \S5 into
explicit morphism data and clarifies the role of the junction spaces.
(\href{https://grothendieck.umontpellier.fr/web/viewer.html?file=../119.pdf\#page=59}{Copy 1, pp.~56--57, beginning at site 58};
\href{https://grothendieck.umontpellier.fr/web/viewer.html?file=../119.pdf\#page=60}{ordered and categorical descriptions on p.~57, site 59};
\href{https://grothendieck.umontpellier.fr/web/viewer.html?file=../119.pdf\#page=37}{call in \S5, p.~34, site 36}.)

\end{document}
